\documentclass[11pt]{article}
\usepackage[letterpaper,margin=1in]{geometry}
\usepackage[T1]{fontenc}
\usepackage[utf8]{inputenc}
\usepackage{lmodern,microtype}
\usepackage{amsmath,amssymb,mathtools,bm,slashed}
\usepackage{xcolor,enumitem,booktabs,array,needspace}
\usepackage[most]{tcolorbox}
\usepackage{fancyhdr,titlesec,parskip,hyperref}
\definecolor{qftblue}{HTML}{245A8D}
\definecolor{qftgreen}{HTML}{2F6B4F}
\definecolor{qftgray}{HTML}{F4F6F8}
\hypersetup{colorlinks=true,linkcolor=qftblue,urlcolor=qftblue,pdfauthor={Nima Lashkari}}
\pagestyle{fancy}\fancyhf{}
\lhead{PHYS 662: Quantum Field Theory I}\cfoot{\thepage}
\setlength{\headheight}{14pt}
\titleformat{\section}{\Large\bfseries\color{qftblue}}{\thesection.}{0.6em}{}
\titleformat{\subsection}{\large\bfseries\color{qftgreen}}{\thesubsection}{0.6em}{}
\setlist[enumerate]{leftmargin=2em,itemsep=0.25em,topsep=0.35em}
\setlist[itemize]{leftmargin=2em,itemsep=0.2em,topsep=0.35em}
\newcommand{\dd}{\mathrm d}
\newcommand{\ii}{\mathrm i}
\newcommand{\eps}{\varepsilon}
\newcommand{\tr}{\operatorname{tr}}
\newcommand{\Tr}{\operatorname{Tr}}
\newcommand{\diag}{\operatorname{diag}}
\newcommand{\Lag}{\mathcal L}
\newcommand{\D}{\mathcal D}
\newcommand{\vac}{\lvert0\rangle}
\newcommand{\vbra}{\langle0\rvert}
\newcommand{\Tord}{\mathrm T}
\newcommand{\ket}[1]{\lvert#1\rangle}
\newcommand{\bra}[1]{\langle#1\rvert}
\newtcolorbox{keyidea}[1][]{enhanced,breakable,colback=qftgray,colframe=qftblue,
 boxrule=0.7pt,arc=1.5mm,title={#1},fonttitle=\bfseries,coltitle=white}
\newtcolorbox{remarkbox}[1][]{enhanced,breakable,colback=white,colframe=qftgreen,
 boxrule=0.6pt,arc=1.5mm,title={#1},fonttitle=\bfseries,coltitle=white}
\newcommand{\module}[1]{\clearpage\section{#1}}
\newcommand{\lecturefront}[4]{%
\rhead{Lecture #1}\hypersetup{pdftitle={PHYS 662 Lecture #1: #2}}
\begin{document}\pagenumbering{gobble}
\begin{titlepage}\centering\vspace*{1.05in}
{\scshape\Large PHYS 662 --- Quantum Field Theory I\par}
\vspace{0.35in}{\Huge\bfseries\color{qftblue}Lecture #1\par}
\vspace{0.2in}{\LARGE\bfseries #2\par}
\vspace{0.2in}{\large\color{qftgreen}#3\par}
\vspace{0.6in}{\large Nima Lashkari\par}{\large Fall 2026\par}
\vfill\begin{minipage}{0.86\textwidth}\small #4\par\medskip
Course page: \url{https://nimalashkari.com/quantum-field-theory-2026-part-i/}
\end{minipage}\vspace{0.5in}\end{titlepage}
\pagenumbering{roman}\setcounter{tocdepth}{1}\tableofcontents
\vfill\begin{remarkbox}[Place in the course]
Lectures 17--18 develop the quantum-mechanical construction. Lecture 19
extends Gaussian integration and source methods to scalar fields.
Lecture 20 introduces Grassmann integration and fermionic fields.
Problem Set 6 brings these four steps together.
\end{remarkbox}\clearpage\pagenumbering{arabic}}
\lecturefront{19}{Path Integrals for the Free Scalar Field}
{Gaussian functionals, Wick contractions, and connected correlations}
{A regulated field is a collection of coupled quantum coordinates. We
derive its path integral, recover the canonical Feynman propagator,
and organize correlations using sources. Euclidean decay then connects
the Gaussian inverse to the mass spectrum.}

\section{From coupled coordinates to a field}
We now set \(\hbar=c=1\), use \(\eta_{\mu\nu}=\diag(1,-1,-1,-1)\),
and keep the propagator convention of Lecture 16:
\begin{equation}
 D_F(x-y)=\langle0|\Tord\phi(x)\phi(y)|0\rangle
 =\int\frac{\dd^4p}{(2\pi)^4}
 \frac{\ii e^{-\ii p\cdot(x-y)}}{p^2-m^2+\ii0}.       \label{eq:DF}
\end{equation}
Thus \((\Box+m^2)D_F=-\ii\delta^{(4)}\). No additional factor of
\(\ii\) is hidden in the definition of \(D_F\).

Start with a finite periodic spatial lattice of spacing \(a\), in
\(d\) spatial dimensions. At site \(\bm n\), set
\(P_{\bm n}=a^d\pi_{\bm n}\); then
\([\phi_{\bm n},P_{\bm m}]=\ii\delta_{\bm n\bm m}\).
The scalar Hamiltonian is
\begin{equation}
 H_a=\sum_{\bm n}\left[\frac{P_{\bm n}^2}{2a^d}
 +\frac{a^d}{2}\left(\sum_{r=1}^d
 \frac{(\phi_{\bm n+\hat r}-\phi_{\bm n})^2}{a^2}
 +m^2\phi_{\bm n}^2\right)\right].
\end{equation}
It is a finite collection of coupled oscillators. Insert field and
momentum resolutions on every time slice exactly as in Lecture 17:
\begin{equation}
 K[\varphi_f,\varphi_i]=\int\D\phi\,\D\pi\,
 \exp\!\left[\ii\int\dd t\dd^d x\,
 (\pi\dot\phi-\mathcal H)\right].
\end{equation}
The boundary labels \(\varphi_i(\bm x),\varphi_f(\bm x)\) are whole
spatial configurations. They are not values of the field at a single
spacetime point.

Each momentum Gaussian produces a factor
\((a^d/(2\pi\ii\epsilon))^{1/2}\) per site and time step. After
integrating momenta, the continuum notation is
\begin{equation}
 K[\varphi_f,\varphi_i]=\int_{\varphi_i}^{\varphi_f}\D\phi\,e^{\ii S_0[\phi]},
 \qquad S_0=\frac12\int\dd^4x\,
 (\partial_\mu\phi\partial^\mu\phi-m^2\phi^2).
                                                        \label{eq:action}
\end{equation}
The lattice measure and the time-slicing limit define \(\D\phi\).
Working at finite regulator makes determinants and integration by
parts ordinary finite-dimensional operations. Absolute vacuum
normalizations may diverge as the regulator is removed, and local
coincident products require separate care even in the free theory.

\module{The vacuum Gaussian and its source}
Project the endpoints onto the vacuum using the contour prescription
of Lecture 18. Integrating \eqref{eq:action} by parts gives a quadratic
operator with Feynman boundary conditions:
\begin{equation}
 S_0+\frac{\ii\varepsilon}{2}\int\phi^2
 =\frac12\phi A\phi,\qquad
 A=-(\Box+m^2)+\ii\varepsilon,\qquad D_F=\ii A^{-1}.
                                                        \label{eq:A}
\end{equation}
The shorthand \(\phi A\phi\) includes spacetime integrations. The
\(+\ii\varepsilon\) in the action yields
\(e^{-\varepsilon\int\phi^2/2}\) in the weight. Boundary wavefunctions
or vacuum projection, not arbitrary fixed endpoints, select this inverse.

For an ordinary real, symmetric positive-definite matrix \(B\), completing the square gives
\begin{equation}
 \int\frac{\dd^n x}{(2\pi)^{n/2}}
 e^{-\frac12x^{\mathsf T}Bx+j^{\mathsf T}x}
 =\frac{e^{\frac12j^{\mathsf T}B^{-1}j}}{\sqrt{\det B}}.
\end{equation}
Analytic continuation to the regulated oscillatory Gaussian replaces
\(B\) by \(-\ii A\). At finite regulator its zero-source integral is
proportional to \(\det(-\ii A)^{-1/2}\), with the phase fixed by the
continuation. The normalization depends on the measure.

Define the normalized vacuum functional
\begin{equation}
 Z_0[J]=\frac{\int\D\phi\,e^{\frac\ii2\phi A\phi+\ii J\phi}}
 {\int\D\phi\,e^{\frac\ii2\phi A\phi}},\qquad Z_0[0]=1.
\end{equation}
The shift \(\phi=\chi-A^{-1}J\) gives
\begin{equation}
 \frac12\phi A\phi+J\phi
 =\frac12\chi A\chi-\frac12JA^{-1}J.
\end{equation}
The determinant cancels between numerator and denominator. Therefore
\begin{keyidea}[Free scalar generating functional]
\begin{equation}
 Z_0[J]=\exp\!\left[-\frac{\ii}{2}JA^{-1}J\right]
       =\exp\!\left[-\frac12\int\dd^4x\dd^4y\,
                      J(x)D_F(x-y)J(y)\right].          \label{eq:Z}
\end{equation}
The two equal expressions differ only in whether the inverse is written
as \(A^{-1}\) or as the physical two-point function \(D_F=\ii A^{-1}\).
\end{keyidea}
The finite-regulator identity is the substantive calculation. The
functional expression is its limit, with sources understood as suitable
test functions so that distributional kernels can act on them.

\module{Correlation functions and Wick's theorem}
Each source derivative inserts \(\ii\phi\):
\begin{equation}
 G_n(x_1,\ldots,x_n)=\langle0|\Tord\phi(x_1)\cdots\phi(x_n)|0\rangle
 =\left.\frac{1}{\ii^n}
 \frac{\delta^n Z_0}{\delta J(x_1)\cdots\delta J(x_n)}\right|_{J=0}.
                                                        \label{eq:moments}
\end{equation}
For two insertions, differentiating \eqref{eq:Z} gives
\(\delta^2 Z_0|_0=-D_F\), and the factor \(1/\ii^2=-1\) restores
the covariance \(G_2=D_F\). The odd functions vanish because
the zero-source Gaussian is centered and invariant under \(\phi\to-\phi\).
They need not vanish in a shifted state or a theory without this symmetry.

For four insertions, write \(D_{ij}=D_F(x_i-x_j)\). Then
\begin{equation}
 G_4=D_{12}D_{34}+D_{13}D_{24}+D_{14}D_{23}.          \label{eq:four}
\end{equation}
For \(2n\) insertions, every term pairs all labels once, giving
\((2n-1)!!=(2n)!/(2^n n!)\) terms. This is Wick's theorem for free
vacuum correlation functions. The combinatorics follow by expanding
the exponential of a quadratic form: each quadratic factor supplies
one pair, and differentiation distributes the external labels.

\subsection{A diagram is a record of a contraction}
Represent each \(D_{ij}\) by a line joining labels \(i\) and \(j\).
The three terms in \eqref{eq:four} are the three pairings of four
labeled endpoints. These are not pictures of three different
classical trajectories; they are a compact notation for products of
distributional kernels. Integration vertices arise only when an
interaction or a composite insertion is introduced.

\begin{center}
\begin{tikzpicture}[x=1cm,y=1cm,every node/.style={font=\small}]
\foreach \shift/\pair in {0/1,4/2,8/3}{
 \begin{scope}[xshift=\shift cm]
 \coordinate (a) at (0,0);\coordinate (b) at (1.5,0);
 \coordinate (c) at (0,-0.85);\coordinate (d) at (1.5,-0.85);
 \ifnum\pair=1
  \draw[qftblue,thick] (a)--(b) (c)--(d);
  \node at (0.75,-1.4) {\(D_{12}D_{34}\)};
 \fi
 \ifnum\pair=2
  \draw[qftblue,thick] (a)--(c) (b)--(d);
  \node at (0.75,-1.4) {\(D_{13}D_{24}\)};
 \fi
 \ifnum\pair=3
  \draw[qftblue,thick] (a)--(d);
  \draw[white,line width=3pt] (b)--(c);
  \draw[qftblue,thick] (b)--(c);
  \node at (0.75,-1.4) {\(D_{14}D_{23}\)};
 \fi
 \foreach \pos/\lab in {a/1,b/2,c/3,d/4}{
  \fill[qftblue] (\pos) circle (1.6pt);
 }
 \node[above] at (a) {1};\node[above] at (b) {2};
 \node[below] at (c) {3};\node[below] at (d) {4};
 \end{scope}
}
\end{tikzpicture}
\end{center}
The crossing in the third pairing is not an interaction vertex.

The operator version also retains normal-ordered terms:
\begin{equation}
 \Tord e^{\ii\int J\hat\phi}
 =:\!e^{\ii\int J\hat\phi}\!:\,
   e^{-\frac12JD_FJ}.                               \label{eq:normal}
\end{equation}
Normal ordering is relative to the free vacuum. Its vacuum expectation
is one, reducing \eqref{eq:normal} to \eqref{eq:Z}. At the operator
level, a time-ordered product is not equal merely to the sum of
fully contracted numbers; uncontracted normal-ordered fields remain.

\begin{remarkbox}[Coincident insertions]
\(D_F\) is a distribution. Expressions such as \(D_F(0)\), or
\(\phi(x)^2\) without a prescription, are not automatically finite.
The regulator, smearing, or a definition such as normal ordering must
be retained when computing composite observables.
\end{remarkbox}

\module{Connected functions and the logarithm}
For any normalized generating functional with \(Z[0]=1\), define
\begin{equation}
 C[J]=\log Z[J],\qquad W[J]=-\ii\log Z[J].
\end{equation}
Near zero source, the connected functions are defined by the cumulant
expansion
\begin{equation}
 C[J]=\sum_{n\geq1}\frac{\ii^n}{n!}
 \int\dd^4x_1\cdots\dd^4x_n\,
 G_{n,c}(x_1,\ldots,x_n)J(x_1)\cdots J(x_n).          \label{eq:cumulants}
\end{equation}
Exponentiating this expansion partitions the labels among connected
groups. For example,
\begin{align}
 G_1(1)&=G_{1,c}(1),\\
 G_2(1,2)&=G_{2,c}(1,2)+G_1(1)G_1(2).
\end{align}
When the one-point function vanishes,
\begin{equation}
 G_4=G_{4,c}+G_{2,c}(1,2)G_{2,c}(3,4)
 +G_{2,c}(1,3)G_{2,c}(2,4)+G_{2,c}(1,4)G_{2,c}(2,3).
\end{equation}
The connected function removes all contributions in which the
insertions split into statistically disconnected groups.

Careful tracking of the factor \(-\ii\) gives
\begin{equation}
 \left.\frac{\delta^n C}{\delta J_1\cdots\delta J_n}\right|_0
 =\ii^nG_{n,c},\qquad
 \left.\frac{\delta^n W}{\delta J_1\cdots\delta J_n}\right|_0
 =\ii^{n-1}G_{n,c}.                                  \label{eq:W}
\end{equation}
For instance, \(\delta^2W|_0=\ii D_F\) in the free vacuum, rather
than \(D_F\) itself. One may choose a different definition of the
connected generator, but its derivative convention must change with it.

For the free real scalar,
\begin{equation}
 C_0[J]=-\tfrac12JD_FJ,\qquad W_0[J]=\tfrac\ii2JD_FJ.
\end{equation}
Only the connected two-point function survives. Equation
\eqref{eq:four} is nevertheless nonzero: its three terms are
disconnected pairings. A Gaussian theory has no connected higher
correlations, not no higher correlations.

\begin{keyidea}[Two different cancellations]
Dividing by the zero-source vacuum functional cancels vacuum diagrams
with no external insertions. Taking a logarithm then selects connected
groups of source insertions. Normalization and connectedness are
related bookkeeping steps, but they are not the same operation.
\end{keyidea}

\module{Field equations inside correlation functions}
The scalar Schwinger--Dyson identity is ordinary integration by parts
at finite regulator. For any insertion \(F[\phi]\),
\begin{equation}
 \int\D\phi\,F\frac{\delta S}{\delta\phi(x)}e^{\ii S}
 =\ii\int\D\phi\,\frac{\delta F}{\delta\phi(x)}e^{\ii S}.
                                                        \label{eq:sd}
\end{equation}
The plus sign on the right is the same as in Lecture 17 with
\(\hbar=1\). For the free action,
\(\delta S_0/\delta\phi=-(\Box+m^2)\phi\).
Taking \(F=\phi(y)\) yields
\begin{equation}
 (\Box_x+m^2)\langle0|\Tord\phi(x)\phi(y)|0\rangle
 =-\ii\delta^{(4)}(x-y),                             \label{eq:greensign}
\end{equation}
which checks both the factor and sign in \eqref{eq:DF}.
The operator equation holds without a source at distinct points;
differentiating time ordering produces the contact term at coincidence.

With \(\mathcal L_{\rm int}=-\lambda\phi^4/4!\), the same argument
would give
\begin{equation}
 (\Box_x+m^2)G_2(x,y)
 +\frac\lambda{3!}\langle\Tord\phi(x)^3\phi(y)\rangle
 =-\ii\delta^{(4)}(x-y).                             \label{eq:hierarchy}
\end{equation}
The regulated composite field must be defined before taking a continuum
limit. This example explains why interacting Schwinger--Dyson equations
form a hierarchy: a two-point equation can involve a four-field
expectation value.

\subsection{The source equation checks the completed square}
Integration by parts with a nonzero source gives
\begin{equation}
 \left[\frac1\ii\int\dd^4y\,A(x,y)\frac{\delta}{\delta J(y)}
       +J(x)\right]Z_0[J]=0.                         \label{eq:sourceSD}
\end{equation}
Since \(\delta Z_0/\delta J=-D_FJZ_0\) and \(AD_F=\ii\), the two
terms cancel. This is a compact independent sign check on the Gaussian
formula. The source-free differential equation alone does not choose
which inverse of \(A\) to use; vacuum boundary conditions still matter.

In a field theory, integration by parts also assumes that the regulator
and measure respect the variation being used. Simple translations of
a free scalar integration variable have unit Jacobian. More general
transformations can act nontrivially on the measure, a question that
will become important when discussing quantum symmetries.

\module{A complex scalar and charged sources}
Write \(\Phi=(\phi_1+\ii\phi_2)/\sqrt2\) with two real fields of
the same mass. Then
\begin{equation}
 S_0[\Phi]=\int\dd^4x\,
 (\partial_\mu\Phi^*\partial^\mu\Phi-m^2\Phi^*\Phi)
 =\Phi^*A\Phi
\end{equation}
with the same Feynman regulator as before. The factor \(1/2\) has
disappeared because one complex field contains two real fields.
At a finite regulator, the two real Gaussian determinants multiply:
\begin{equation}
 \mathcal Z_{0,\mathbb C}[0]=\mathcal N\,[\det(-\ii A)]^{-1}.
\end{equation}
The constant \(\mathcal N\) depends on the measure convention. This
unnormalized vacuum integral is not the normalized value \(Z[0]=1\).

Introduce independent commuting sources \(\bar J,J\), differentiating
them independently even when they will eventually be complex conjugates:
\begin{equation}
 Z_0[\bar J,J]=
 \frac{\int\D\Phi^*\D\Phi\,
 e^{\ii(\Phi^*A\Phi+\bar J\Phi+\Phi^*J)}}{\mathcal Z_{0,\mathbb C}[0]}.
\end{equation}
The shift \(\Phi=\chi-A^{-1}J\),
\(\Phi^*=\chi^*-\bar JA^{-1}\) gives
\begin{equation}
 Z_0[\bar J,J]=e^{-\ii\bar JA^{-1}J}=e^{-\bar JD_FJ}.
                                                        \label{eq:complexZ}
\end{equation}
One derivative with respect to \(\bar J(x)\) inserts \(\ii\Phi(x)\),
and one with respect to \(J(y)\) inserts \(\ii\Phi^*(y)\). Hence
\begin{align}
 \langle\Tord\Phi(x)\Phi^\dagger(y)\rangle&=D_F(x-y),\\
 \langle\Tord\Phi(x)\Phi(y)\rangle
 =\langle\Tord\Phi^\dagger(x)\Phi^\dagger(y)\rangle&=0.
\end{align}
The vanishing charged pairings follow from the unbroken global
\(U(1)\) symmetry of the vacuum. They can also be checked by expanding
in \(\phi_1,\phi_2\).

For example,
\begin{equation}
 \langle\Tord\Phi_1\Phi_2\Phi^\dagger_3\Phi^\dagger_4\rangle
 =D_{13}D_{24}+D_{14}D_{23}.                          \label{eq:complex4}
\end{equation}
There are two allowed pairings, not three: charge flow excludes a
\(\Phi\Phi\) contraction. This is a useful bridge to Dirac fields,
where sources also come in two types but are Grassmann odd, changing
the relative sign between the two pairings.

\module{The Euclidean scalar and its inverse}
Under \(t=-\ii\tau\), the free Euclidean action is
\begin{equation}
 S_E=\frac12\int\dd^Dx_E\,
 [ (\partial_E\phi)^2+m^2\phi^2]
 =\frac12\phi A_E\phi,\qquad A_E=-\partial_E^2+m^2,
\end{equation}
where \(D=d+1\) is the total Euclidean dimension. For \(m>0\) the
finite-volume quadratic form is strictly positive. Its covariance is
\begin{equation}
 C_E(x)=A_E^{-1}(x)=\int\frac{\dd^Dp}{(2\pi)^D}
 \frac{e^{\ii p\cdot x}}{p^2+m^2},\qquad A_EC_E=\delta^{(D)}.
                                                        \label{eq:CE}
\end{equation}
The normalized real-source functional has a positive quadratic exponent:
\begin{equation}
 Z_E[J]=\frac{\int\D\phi\,e^{-S_E+J\phi}}{\int\D\phi\,e^{-S_E}}
 =e^{\frac12JC_EJ}.
\end{equation}
Its ordinary derivatives insert \(\phi\), with no factors of \(\ii\).
The sign change relative to Minkowski sources is therefore expected.
Analytic continuation includes the energy contour and its measure;
substituting \(p^0=\ii p_D\) in a denominator alone misses a factor.

\subsection{Proper time turns a propagator into a Gaussian}
For separated points \(R=|x|>0\), use
\(1/(p^2+m^2)=\int_0^\infty\dd s\,e^{-s(p^2+m^2)}\).
The momentum integral then yields
\begin{equation}
 C_E(R)=\int_0^\infty\frac{\dd s}{(4\pi s)^{D/2}}
 \exp\!\left(-m^2s-\frac{R^2}{4s}\right).             \label{eq:propertime}
\end{equation}
This representation displays ultraviolet sensitivity at small \(s\)
and long-distance damping at large \(R\). At coincidence, the
small-\(s\) integral diverges for \(D\geq2\), even though the
Gaussian theory is exactly solvable.

At temperature \(\beta^{-1}\), bosonic fields are periodic in
Euclidean time with circumference \(\beta\). Replace the energy
integral by \(\beta^{-1}\sum_{n\in\mathbb Z}\), with
\(p_D=2\pi n/\beta\). At finite volume, the massless constant mode
has zero action and needs an infrared prescription; one cannot
blindly invert it. On the infinite line at zero temperature,
\eqref{eq:CE} is the continuation of the vacuum Feynman function.

\module{Bessel functions, distance scales, and the mass gap}
Evaluating \eqref{eq:propertime} gives
\begin{equation}
 C_E(R)=\frac{m^\nu}{(2\pi)^{\nu+1}R^\nu}K_\nu(mR),
 \qquad \nu=\frac D2-1=\frac{d-1}{2}.                \label{eq:bessel}
\end{equation}
The order is \((d-1)/2\) when \(d\) denotes spatial dimension.
For \(D=4\), this becomes \(mK_1(mR)/(4\pi^2R)\).
Using \(K_\nu(z)\sim\sqrt{\pi/(2z)}e^{-z}\) at large \(z\),
\begin{equation}
 C_E(R)\sim
 \frac{m^{(D-3)/2}}{2(2\pi)^{(D-1)/2}}
 \frac{e^{-mR}}{R^{(D-1)/2}}.
                                                        \label{eq:largeR}
\end{equation}
The power is dimension dependent; a universal \(e^{-mR}/R\)
would be incorrect. At short distance, for \(D>2\),
\begin{equation}
 C_E(R)\sim\frac{\Gamma(D/2-1)}{4\pi^{D/2}R^{D-2}}.
\end{equation}
In \(D=2\), \(C_E=K_0(mR)/(2\pi)\) instead behaves as
\(-[\log(mR/2)+\gamma_E]/(2\pi)\). This logarithm is also why
the massless limit needs an infrared convention in two dimensions.

\subsection{Reading physical information from decay}
For a general vacuum, an operator's connected Euclidean correlator
has a spectral expansion in energy differences. Large separation
selects the lightest state with a nonzero matrix element in that
channel. The elementary free scalar creates a one-particle state of
mass \(m\), so its correlation length is \(1/m\). A normal-ordered
bilinear instead couples first to two-particle states and can decay
with threshold \(2m\). A theory's mass gap and an individual
operator's leading decay scale must therefore be distinguished.

\subsection*{Questions to check understanding}
\begin{enumerate}
\item Which lattice canonical momentum has commutator \(\ii\delta_{\bm n\bm m}\)?
\item Why does a real Gaussian give a determinant to power \(-1/2\),
whereas a complex Gaussian gives power \(-1\)?
\item Verify the source and Schwinger--Dyson checks of \(D_F\).
\item Why does \(\log Z\) generate connected functions rather than all functions?
\item Compare the three real-scalar pairings with the two charged pairings.
\item What sets the exponential and the power law in \eqref{eq:largeR}?
\end{enumerate}
Next lecture replaces commuting integration variables by anticommuting
ones, while retaining this source-and-inverse strategy.
\end{document}
